% GENERATED FILE. Edit canonical lessons, then run npm run generate:books.
% canonical-source-hash: fnv1a64:ae53555add4c98ec
% canonical-lessons: TE-C445-rahadari, TE-C445-lari, TE-C445-dusukupovu, TE-C445-nemmadincu, TE-C445-adupu-tappu, TE-R445-review-438-445

\chapter{On the Highway}
\label{ch:te-on-the-highway}
\hlchaptermodality{445}

\begin{chapteropening}
\textbf{By the end of this chapter:} \emph{I can say a highway, a lorry, to shoot ahead, to slow down and to lose control.}

Complete the last lesson of chapter 445: I can say a highway, a lorry, to shoot ahead, to slow down and to lose control.
\end{chapteropening}

\section[\texorpdfstring{\te{రహదారి} (rahadāri)}{రహదారి (rahadāri)}]{\te{రహదారి} (rahadāri) — a highway; a main road}

\label{lesson:TE-C445-rahadari}



\begin{quote}
Before the new one: say the Telugu for a puncture, then the Telugu for to mend.
\end{quote}

\subsection*{You'll want to know: \te{రహదారి}}
\textbf{\te{రహదారి}} — \emph{rahadāri} — \textquotedblleft{}a highway; a main road\textquotedblright{}.

\begin{quote}
\te{జాతీయ} \te{రహదారి} \te{మీద} \te{వాహనాలు} \te{వేగంగా} \te{వెళ్తాయి}.
\end{quote}

\begin{quote}
\emph{jātīya rahadāri mīda vāhanālu vēgaṅgā veḷtāyi.}
\end{quote}

\begin{quote}
Vehicles go fast on the national highway.
\end{quote}

\te{జాతీయ} \te{రహదారి} is a national highway. \te{దారి} alone is any way or path.

\begin{cousinweb}
\textbf{\te{రహదారి}} (\emph{rahadāri}) is a compound whose second part is \te{దారి}, \textquotedblleft{}a way, a road\textquotedblright{}; the first part is usually traced to Persian \emph{rāh}, \textquotedblleft{}a road\textquotedblright{}, and the whole is read as \textquotedblleft{}the main road\textquotedblright{}.
\end{cousinweb}

\begin{grammarlens}[title={What you've built}]
One more word for getting around.
\end{grammarlens}

\subsection*{Your turn}
\begin{itemize}
\raggedright
  \item \emph{Say it:} \emph{rahadāri}
  \item \emph{Say it:} \emph{rahadāri}, once more
  \item \emph{Say it:} \emph{jātīya rahadāri mīda vāhanālu vēgaṅgā veḷtāyi.} — Vehicles go fast on the national highway.
  \item \emph{Recall:} say the Telugu for a puncture, then the Telugu for to mend, then say \emph{rahadāri} again
\end{itemize}

\begin{tcolorbox}[breakable,colback=teal!4,colframe=teal!35!black,title={Before you move on}]
What does \te{రహదారి} mean? (\textquotedblleft{}A highway; a main road\textquotedblright{}.) Say it once more.
\end{tcolorbox}
\section[\texorpdfstring{\te{లారీ} (lārī)}{లారీ (lārī)}]{\te{లారీ} (lārī) — a lorry; a truck}

\label{lesson:TE-C445-lari}



\begin{quote}
Before the new one: say the Telugu for to mend, then the Telugu for a highway.
\end{quote}

\subsection*{You'll want to know: \te{లారీ}}
\textbf{\te{లారీ}} — \emph{lārī} — \textquotedblleft{}a lorry; a truck\textquotedblright{}.

\begin{quote}
\te{లారీ} \te{నిండా} \te{ఇసుక} \te{ఉంది}.
\end{quote}

\begin{quote}
\emph{lārī niṇḍā isuka undi.}
\end{quote}

\begin{quote}
The lorry is full of sand.
\end{quote}

Plural: \te{లారీలు}, \emph{lārīlu}.

\te{నిండా}, full of, follows the load: \te{ఇసుక}, sand.

\begin{cousinweb}
\textbf{\te{లారీ}} (\emph{lārī}) is a loan from English \emph{lorry}, a goods truck.
\end{cousinweb}

\begin{grammarlens}[title={What you've built}]
One more word for getting around.
\end{grammarlens}

\subsection*{Your turn}
\begin{itemize}
\raggedright
  \item \emph{Say it:} \emph{lārī}
  \item \emph{Say it:} \emph{lārī}, once more
  \item \emph{Say it:} \emph{lārī niṇḍā isuka undi.} — The lorry is full of sand.
  \item \emph{Recall:} say the Telugu for to mend, then the Telugu for a highway, then say \emph{lārī} again
\end{itemize}

\begin{tcolorbox}[breakable,colback=teal!4,colframe=teal!35!black,title={Before you move on}]
What does \te{లారీ} mean? (\textquotedblleft{}A lorry; a truck\textquotedblright{}.) Say it once more.
\end{tcolorbox}
\section[\texorpdfstring{\te{దూసుకుపోవు} (dūsukupōvu)}{దూసుకుపోవు (dūsukupōvu)}]{\te{దూసుకుపోవు} (dūsukupōvu) — to shoot ahead; to dash past}

\label{lesson:TE-C445-dusukupovu}



\begin{quote}
Before the new one: say the Telugu for a highway, then the Telugu for a lorry.
\end{quote}

\subsection*{You'll want to know: \te{దూసుకుపోవు}}
\textbf{\te{దూసుకుపోవు}} — \emph{dūsukupōvu} — \textquotedblleft{}to shoot ahead; to dash past\textquotedblright{}.

\begin{quote}
\te{ఒక} \te{కారు} \te{మా} \te{బస్సును} \te{దాటి} \te{దూసుకుపోయింది}.
\end{quote}

\begin{quote}
\emph{oka kāru mā bassunu dāṭi dūsukupōyindi.}
\end{quote}

\begin{quote}
A car overtook our bus and shot ahead.
\end{quote}

Past (I …): \te{దూసుకుపోయాను}, \emph{dūsukupōyānu}.

Said of something moving fast and pushing through: a car, a ball, a crowd.

\begin{cousinweb}
\textbf{\te{దూసుకుపోవు}} (\emph{dūsukupōvu}) is a Telugu compound: its first part has an origin that is not settled, and its second part is the native \te{పోవు}, \textquotedblleft{}to go\textquotedblright{} (DEDR 4572).
\end{cousinweb}

\begin{grammarlens}[title={What you've built}]
One more word for getting around.
\end{grammarlens}

\subsection*{Your turn}
\begin{itemize}
\raggedright
  \item \emph{Say it:} \emph{dūsukupōvu}
  \item \emph{Say it:} \emph{dūsukupōvu}, once more
  \item \emph{Say it:} \emph{oka kāru mā bassunu dāṭi dūsukupōyindi.} — A car overtook our bus and shot ahead.
  \item \emph{Recall:} say the Telugu for a highway, then the Telugu for a lorry, then say \emph{dūsukupōvu} again
\end{itemize}

\begin{tcolorbox}[breakable,colback=teal!4,colframe=teal!35!black,title={Before you move on}]
What does \te{దూసుకుపోవు} mean? (\textquotedblleft{}To shoot ahead; to dash past\textquotedblright{}.) Say it once more.
\end{tcolorbox}
\section[\texorpdfstring{\te{నెమ్మదించు} (nemmadiñcu)}{నెమ్మదించు (nemmadiñcu)}]{\te{నెమ్మదించు} (nemmadiñcu) — to slow down}

\label{lesson:TE-C445-nemmadincu}



\begin{quote}
Before the new one: say the Telugu for a lorry, then the Telugu for to shoot ahead.
\end{quote}

\subsection*{You'll want to know: \te{నెమ్మదించు}}
\textbf{\te{నెమ్మదించు}} — \emph{nemmadiñcu} — \textquotedblleft{}to slow down\textquotedblright{}.

\begin{quote}
\te{మలుపు} \te{దగ్గర} \te{బస్సు} \te{నెమ్మదించింది}.
\end{quote}

\begin{quote}
\emph{malupu daggara bassu nemmadiñcindi.}
\end{quote}

\begin{quote}
At the bend the bus slowed down.
\end{quote}

Past (I …): \te{నెమ్మదించాను}, \emph{nemmadiñcānu}.

\te{నెమ్మదిగా}, slowly, is the same word. It is also said of rain that eases and of prices that stop rising.

\begin{cousinweb}
\textbf{\te{నెమ్మదించు}} (\emph{nemmadiñcu}) is native Dravidian: \te{నెమ్మది}, \textquotedblleft{}tranquillity, ease\textquotedblright{} (DEDR 3602), made into a verb with -\te{ఇంచు}.
\end{cousinweb}

\begin{grammarlens}[title={What you've built}]
One more word for getting around.
\end{grammarlens}

\subsection*{Your turn}
\begin{itemize}
\raggedright
  \item \emph{Say it:} \emph{nemmadiñcu}
  \item \emph{Say it:} \emph{nemmadiñcu}, once more
  \item \emph{Say it:} \emph{malupu daggara bassu nemmadiñcindi.} — At the bend the bus slowed down.
  \item \emph{Recall:} say the Telugu for a lorry, then the Telugu for to shoot ahead, then say \emph{nemmadiñcu} again
\end{itemize}

\begin{tcolorbox}[breakable,colback=teal!4,colframe=teal!35!black,title={Before you move on}]
What does \te{నెమ్మదించు} mean? (\textquotedblleft{}To slow down\textquotedblright{}.) Say it once more.
\end{tcolorbox}
\section[\texorpdfstring{\te{అదుపు} \te{తప్పు} (adupu tappu)}{అదుపు తప్పు (adupu tappu)}]{\te{అదుపు} \te{తప్పు} (adupu tappu) — to lose control}

\label{lesson:TE-C445-adupu-tappu}



\begin{quote}
Before the new one: say the Telugu for to shoot ahead, then the Telugu for to slow down.
\end{quote}

\subsection*{You'll want to know: \te{అదుపు} \te{తప్పు}}
\textbf{\te{అదుపు} \te{తప్పు}} — \emph{adupu tappu} — \textquotedblleft{}to lose control\textquotedblright{}.

\begin{quote}
\te{వర్షంలో} \te{కారు} \te{అదుపు} \te{తప్పింది}.
\end{quote}

\begin{quote}
\emph{varṣaṁlō kāru adupu tappindi.}
\end{quote}

\begin{quote}
In the rain the car went out of control.
\end{quote}

Past (I …): \te{అదుపు} \te{తప్పాను}, \emph{adupu tappānu}.

\te{అదుపు}, control, + \te{తప్పు}, to slip. It is said of a vehicle, a fire, a crowd or a temper.

\begin{cousinweb}
\textbf{\te{అదుపు} \te{తప్పు}} (\emph{adupu tappu}) is a Telugu phrase: \te{అదుపు}, \textquotedblleft{}control, restraint\textquotedblright{}, whose origin is not settled, and the native \te{తప్పు}, \textquotedblleft{}to slip, to miss\textquotedblright{} (DEDR 3071).
\end{cousinweb}

\begin{grammarlens}[title={What you've built}]
One more word for getting around.
\end{grammarlens}

\subsection*{Your turn}
\begin{itemize}
\raggedright
  \item \emph{Say it:} \emph{adupu tappu}
  \item \emph{Say it:} \emph{adupu tappu}, once more
  \item \emph{Say it:} \emph{varṣaṁlō kāru adupu tappindi.} — In the rain the car went out of control.
  \item \emph{Recall:} say the Telugu for to shoot ahead, then the Telugu for to slow down, then say \emph{adupu tappu} again
\end{itemize}

\begin{tcolorbox}[breakable,colback=teal!4,colframe=teal!35!black,title={Before you move on}]
What does \te{అదుపు} \te{తప్పు} mean? (\textquotedblleft{}To lose control\textquotedblright{}.) Say it once more.
\end{tcolorbox}
\section[(dialogue)]{Getting around, first steps}

\label{lesson:TE-R445-review-438-445}



\begin{quote}
Say the Telugu for a passenger and to travel.
\end{quote}

\subsection*{Your turn}
Say these aloud:

\begin{itemize}
\raggedright
  \item a passenger, to travel, a destination, a route, to set foot in
  \item a reservation, a booking, a sleeper berth, a railway coach, a fare
  \item a delay, to wait, punctuality, an interval, to hurry up
  \item a passport, identification, a customs duty, security, a certificate
  \item a rail, the metro, comings and goings, to be packed, to derail
  \item a vehicle, an engine, a tyre, a brake, a horn
  \item fuel, to break down, to refuse to start, a puncture, to mend
  \item a highway, a lorry, to shoot ahead, to slow down, to lose control
\end{itemize}

\begin{tcolorbox}[breakable,colback=teal!4,colframe=teal!35!black,title={Before you move on}]
A passenger? (\textbf{\te{ప్రయాణికుడు}}, \emph{prayāṇikuḍu}.) To travel? (\textbf{\te{ప్రయాణించు}}, \emph{prayāṇiñcu}.) A destination? (\textbf{\te{గమ్యం}}, \emph{gamyaṁ}.) A route? (\textbf{\te{మార్గం}}, \emph{mārgaṁ}.) To set foot in? (\textbf{\te{అడుగుపెట్టు}}, \emph{aḍugupeṭṭu}.) A reservation? (\textbf{\te{రిజర్వేషను}}, \emph{rijarvēṣanu}.) A booking? (\textbf{\te{బుకింగు}}, \emph{bukiṅgu}.) A sleeper berth? (\textbf{\te{బెర్తు}}, \emph{bertu}.) A railway coach? (\textbf{\te{బోగీ}}, \emph{bōgī}.) A fare? (\textbf{\te{కిరాయి}}, \emph{kirāyi}.) A delay? (\textbf{\te{ఆలస్యం}}, \emph{ālasyaṁ}.) To wait? (\textbf{\te{వేచి} \te{ఉండు}}, \emph{vēci uṇḍu}.) Punctuality? (\textbf{\te{సమయపాలన}}, \emph{samayapālana}.) An interval? (\textbf{\te{విరామం}}, \emph{virāmaṁ}.) To hurry up? (\textbf{\te{త్వరపడు}}, \emph{tvarapaḍu}.) A passport? (\textbf{\te{పాస్పోర్టు}}, \emph{pāspōrṭu}.) Identification? (\textbf{\te{గుర్తింపు}}, \emph{gurtimpu}.) A customs duty? (\textbf{\te{సుంకం}}, \emph{suṅkaṁ}.) Security? (\textbf{\te{భద్రత}}, \emph{bhadrata}.) A certificate? (\textbf{\te{ధ్రువపత్రం}}, \emph{dhruvapatraṁ}.) A rail? (\textbf{\te{పట్టా}}, \emph{paṭṭā}.) The metro? (\textbf{\te{మెట్రో}}, \emph{meṭrō}.) Comings and goings? (\textbf{\te{రాకపోకలు}}, \emph{rākapōkalu}.) To be packed? (\textbf{\te{కిక్కిరిసిపోవు}}, \emph{kikkirisipōvu}.) To derail? (\textbf{\te{పట్టాలు} \te{తప్పు}}, \emph{paṭṭālu tappu}.) A vehicle? (\textbf{\te{వాహనం}}, \emph{vāhanaṁ}.) An engine? (\textbf{\te{ఇంజను}}, \emph{iñjanu}.) A tyre? (\textbf{\te{టైరు}}, \emph{ṭairu}.) A brake? (\textbf{\te{బ్రేకు}}, \emph{brēku}.) A horn? (\textbf{\te{హారను}}, \emph{hāranu}.) Fuel? (\textbf{\te{ఇంధనం}}, \emph{indhanaṁ}.) To break down? (\textbf{\te{చెడిపోవు}}, \emph{ceḍipōvu}.) To refuse to start? (\textbf{\te{మొరాయించు}}, \emph{morāyiñcu}.) A puncture? (\textbf{\te{పంక్చరు}}, \emph{paṅkcaru}.) To mend? (\textbf{\te{మరమ్మతు} \te{చేయు}}, \emph{marammatu cēyu}.) A highway? (\textbf{\te{రహదారి}}, \emph{rahadāri}.) A lorry? (\textbf{\te{లారీ}}, \emph{lārī}.) To shoot ahead? (\textbf{\te{దూసుకుపోవు}}, \emph{dūsukupōvu}.) To slow down? (\textbf{\te{నెమ్మదించు}}, \emph{nemmadiñcu}.) To lose control? (\textbf{\te{అదుపు} \te{తప్పు}}, \emph{adupu tappu}.)
\end{tcolorbox}
